%======================================================================
%  BACHELOR THESIS PROPOSAL — FOUR-PAGE VERSION
%  Ruhr-Universität Bochum — IT-Sicherheit / Informationstechnik
%
%  This is the version to hand in: the supervisor asked for four pages.
%  `proposal.tex` is the long version and stays in the repository as the
%  full methodological record — every decision compressed here is argued
%  in full there, and the thesis methodology chapter is written from it.
%
%  HARD CONSTRAINT: this document must stay at four pages, references
%  included. After any edit, recompile and check. If it spills onto a
%  fifth page, cut content — do not shrink the margins or the font.
%
%  COMPILE:  pdflatex -> biber -> pdflatex -> pdflatex
%======================================================================
\documentclass[10pt,a4paper]{article}

%----------------------------------------------------------------------
%  Metadata — keep in sync with proposal.tex
%----------------------------------------------------------------------
\newcommand{\thesistitle}{Understanding Privacy Harm Communication in Fitness and Health Tracking Apps}
\newcommand{\authorname}{Hamza}
\newcommand{\matriculation}{TODO: matriculation number}
\newcommand{\firstexaminer}{TODO: Erstprüfer*in (RUB)}
\newcommand{\supervisorname}{Kavous Salehzadeh Niksirat, MPI-SP}
\newcommand{\theprogram}{B.Sc. IT-Sicherheit / Informationstechnik}
\newcommand{\thesemester}{Winter Semester 2026/27}

%----------------------------------------------------------------------
%  Packages
%----------------------------------------------------------------------
\usepackage[utf8]{inputenc}
\usepackage[T1]{fontenc}
\usepackage[english]{babel}
\usepackage{lmodern}
\usepackage{microtype}

\usepackage[a4paper,margin=2.2cm,top=1.8cm,bottom=1.9cm]{geometry}
\usepackage{booktabs}
\usepackage{tabularx}
\usepackage{enumitem}
\usepackage{csquotes}
\usepackage{xcolor}
\usepackage{tikz}
\usetikzlibrary{arrows.meta}

% Tight headings: this is what buys the fourth page back.
\usepackage[compact]{titlesec}
\titlespacing*{\section}{0pt}{1.1ex plus .2ex}{0.5ex}
\titlespacing*{\subsection}{0pt}{0.9ex plus .2ex}{0.4ex}
\titleformat{\section}{\normalfont\large\bfseries}{\thesection}{0.6em}{}
\titleformat{\subsection}{\normalfont\bfseries}{\thesubsection}{0.6em}{}

\setlist{nosep,leftmargin=*}

\definecolor{linkblue}{HTML}{16537E}
\usepackage[hidelinks]{hyperref}
\hypersetup{
  colorlinks = true,
  linkcolor  = linkblue,
  citecolor  = linkblue,
  urlcolor   = linkblue,
  pdftitle   = {\thesistitle},
  pdfauthor  = {\authorname}
}

%----------------------------------------------------------------------
%  Bibliography — same shared file as every other document here.
%  maxbibnames=2 collapses long author lists to "et al." This version
%  needs the space; the long proposal prints them in full.
%----------------------------------------------------------------------
\usepackage[backend=biber,style=numeric-comp,sorting=none,
            maxbibnames=2,minbibnames=1,maxcitenames=1,
            urldate=long]{biblatex}
\addbibresource{../../../_shared/references.bib}
\renewcommand*{\bibfont}{\footnotesize}
\setlength{\bibitemsep}{0pt}

\pagestyle{empty}

%======================================================================
\begin{document}

%---------------------------- Header ----------------------------------
\begin{center}
  {\large\bfseries Understanding Privacy Harm Communication\\
   in Fitness and Health Tracking Apps\par}
  \vspace{0.35em}
  {\normalsize Bachelor's Thesis Proposal\par}
  \vspace{0.5em}
  {\footnotesize
   \authorname\ (\matriculation) \quad·\quad \theprogram \quad·\quad \thesemester\\
   Supervisor: \supervisorname \quad·\quad First examiner: \firstexaminer
   \quad·\quad \today\par}
\end{center}
\vspace{0.4em}
\hrule
\vspace{0.9em}

%---------------------------- 1. Background ---------------------------
\section{Background}

In Germany, 73\,\% of smartphone users report using health apps for fitness,
nutrition, or general wellbeing~\parencite{bitkom2025health}. Fitness, sleep,
and menstrual trackers accumulate three kinds of data at once: attributes users
enter directly (weight, symptoms, mood, sexual activity), continuous sensor
streams (steps, heart rate, sleep, location), and system context (identifiers,
timestamps, usage patterns). Under Article~9(1) GDPR, data concerning health is
a special category of personal data, and menstrual and fertility data falls
squarely within it~\parencite{gdpr2016}.

A privacy \emph{violation} and a privacy \emph{harm} are not the same thing.
Citron and Solove classify the injuries themselves into seven types --- physical,
economic, reputational, psychological, autonomy, discrimination, and
relationship harm~\parencite[830--861]{citron2021harms}. Knabenhans et al.\
supply what a typology alone does not: the route from a system property to an
injury. Their \emph{harm trees} link system components to adversarial
capabilities --- distinguishing between users, tracking them, inferring
attributes about them --- and those capabilities to harms, and they show that
privacy-enhancing technologies (PETs) close only the pathways that come from
implementation choices, never those inherent to the
functionality~\parencite{knabenhans2026harms}. The inference capability is not
hypothetical for trackers: a survey of 236 studies reports that tracker data can
be mined to infer private attributes, personality traits among
them~\parencite{niksirat2024wearable}.

That matters because modern risk is inference-centric rather than
collection-centric. Martin et al.\ examine the three principles behind industry
PET deployments --- limiting third-party access, preferring inferences over raw
data, and keeping data on the device --- and show that all three can be satisfied
while the privacy-violating behaviour continues untouched, because the sensitive
object is the inference, not the raw record~\parencite{martin2025cookies}.

Providers communicate privacy on four surfaces: the privacy policy, the store
listing and its structured disclosure, the provider website and its advertising,
and the app interface itself. The law shapes only the first, and shapes it
narrowly. Articles 13--14 GDPR enumerate the controller, purposes, legal basis,
recipients, transfers, retention, and rights. Consequences appear in exactly one
place --- Article 13(2)(f), for automated decision-making, requires
\enquote{meaningful information about \dots\ the significance and the envisaged
consequences of such processing for the data subject}. Breaches are handled
reactively under Article 34, after the fact~\parencite{gdpr2016}. \textbf{Nothing
obliges a provider to say what can be derived from the data it collects, or what
would follow for the user if that data were exposed.}

%---------------------------- 2. Research Problem ---------------------
\section{Research Problem}

A substantial literature already asks whether providers tell the truth.
Huckvale et al.\ found 29 of 36 health apps transmitting data to Google or
Facebook services while only 12 disclosed it~\parencite{huckvale2019assessment};
a dynamic study of nearly 5{,}000 Android apps found 81\,\% misrepresenting
their practices in the Data Safety
section~\parencite{arkalakis2024abandon}; and menstrual and fertility trackers
specifically have been examined through policy analysis, usability inspection,
and network
measurement~\parencite{shipp2020period,malki2024exploring,hudig2025intimate}.
A second literature asks what users believe: fitness tracker users are confident
about \emph{what} is collected but unsure \emph{how} it is used, and rate
harmful scenarios as plausible yet unlikely~\parencite{gabriele2020understanding};
period tracker users report high concern combined with feeling uninformed and
powerless~\parencite{cao2024deleted}.

Neither literature asks which \emph{harms} the communication names. That is the
gap: \textbf{we lack a systematic account of what harms service providers
communicate} --- which types are acknowledged, which are minimised, which never
appear, and how PETs are framed (secrecy versus harm prevention). Three things
in particular sit outside the frame the law and the platform schemas impose ---
\textbf{inference} (what is derived beyond what the user provided),
\textbf{breach consequences} (what would happen to the user), and
\textbf{incident handling} --- together with a fourth element cutting across
them: whether any \textbf{concrete consequence} is ever named, or whether the
text stays at the level of data flows. Where a provider is silent on all four,
the user is left with material that describes a mechanism and never describes a
stake, which is exactly the communication environment the user studies above
would predict.

\subsection*{Research questions}

\noindent\textbf{RQ.} \emph{What types of downstream privacy harms do health
tracker service providers explicitly acknowledge or downplay when describing
their privacy practices or PETs, and how do they frame the protection these
technologies offer?}

\begin{description}[leftmargin=1.6cm,style=nextline,itemsep=1pt,topsep=2pt]
  \item[\textbf{RQ1}] \textbf{Coverage.} Which harm types appear, and which are
    absent? How often is inference, breach, incident handling, or a concrete
    consequence named at all?
  \item[\textbf{RQ2}] \textbf{Stance.} Where a harm is mentioned, is it
    acknowledged, hedged, minimised, or deflected onto the user?
  \item[\textbf{RQ3}] \textbf{PET framing.} Which of Martin et al.'s three
    principles do providers invoke, and is protection framed as secrecy or as
    harm prevention?
  \item[\textbf{RQ4}] \textbf{Consistency.} Does the answer change between the
    policy, the store listing, the website, and the interface?
\end{description}

\subsection*{Scope}

In scope: fitness and activity tracking (sleep, heart rate, and nutrition
included) and fertility\,/\,period tracking, on the German storefronts of both
stores, analysed through the app rather than the device. Deliberately out of
scope: \emph{a user study} --- the scenario-based survey in the original thesis
offer is a possible second phase, for which this corpus would supply realistic
stimulus material; \emph{verifying whether claims are true}, which traffic and
code analysis already address; \emph{grading legal compliance}, which is a
judgement for lawyers, not a coding task; and \emph{prescription DiGA and
regulated medical devices}, whose communication follows a different regime.

\subsection*{Why it matters}

\textbf{For users:} deciding whether to install a cycle tracker means weighing a
benefit you can feel against a risk you have to imagine. If the communication
never names a consequence, the user must supply it themselves --- and the
evidence says they cannot. A user who reads that an app does not share raw data
may reasonably conclude they are protected; if the provider derives and acts on
an inference instead, that conclusion is wrong and nothing in the text said so.
\textbf{For providers and designers:} a map of which harms are named, on which
surface, in what terms, is an account of what their communication omits,
expressed in categories a non-lawyer can act on. \textbf{For regulators:}
whether harm-facing disclosure would add anything is currently an untested
assumption. If the four elements above are near-absent across a market segment
handling Article~9 data, that is a concrete finding about the limits of a
transparency regime built on data-practice disclosure. \textbf{For research:}
the harm typologies were built for courts and for system analysis; neither has
been turned on the text providers actually put in front of users.

%---------------------------- 3. Methodology --------------------------
\section{Methodology}

A structured content analysis of privacy communication, combined with a
task-based inspection of the interfaces that carry part of it. Every code is
anchored to a verbatim quote with source, location, and capture date; every
instrument --- sampling protocol, codebook, each LLM prompt --- is versioned in
the repository.

\begin{figure}[h]
\centering
\begin{tikzpicture}[
  font=\scriptsize,
  bx/.style ={draw, rounded corners=2pt, align=center, text width=27mm,
              minimum height=16mm, inner sep=2.5pt},
  arr/.style={-{Stealth[length=1.6mm]}, semithick}
]
  \hyphenpenalty=10000 \exhyphenpenalty=10000
  \node[bx] (a) at (-6.6,0) {\textbf{Sample}\\50 apps,\\German stores,\\
                             stratified 30/20};
  \node[bx] (b) at (-3.3,0) {\textbf{Corpus}\\policies, listings,\\sites, ads,
                             plus\\12 walkthroughs};
  \node[bx] (c) at ( 0.0,0) {\textbf{Coding}\\harm type\\$\cdot$ stance
                             $\cdot$\\PET framing};
  \node[bx] (d) at ( 3.3,0) {\textbf{Validation}\\$\kappa$ with prevalence,\\
                             precision \emph{and}\\recall};
  \node[bx] (e) at ( 6.6,0) {\textbf{Answers}\\RQ1--RQ4, every\\code backed
                             by\\a verbatim quote};
  \draw[arr] (a) -- (b); \draw[arr] (b) -- (c);
  \draw[arr] (c) -- (d); \draw[arr] (d) -- (e);
\end{tikzpicture}
\end{figure}

\textbf{Step 1 --- Sampling.} Google Play and the Apple App Store, German
storefronts. A frozen keyword list in German and English covers both domains
(\emph{Schrittzähler, Schlaf-Tracker, Kalorienzähler\dots}; \emph{Periode,
Zyklus, Eisprung, Kinderwunsch\dots}). \enquote{Highest rated} is fixed in
advance: the first 20 results per keyword are recorded, the pooled set filtered
to apps with $\geq 1{,}000$ ratings --- below that the average is noise --- then
ranked by average rating with rating count as tie-break. Target: \textbf{50
apps}, stratified 30 fitness / 20 fertility, the fertility stratum deliberately
over-represented because it is where the harm literature is most developed. A
five-app pilot tests the instruments first. Apps without a reachable policy,
prescription DiGA, medical devices, and apps requiring identity verification are
excluded.

\textbf{Step 2 --- Material.} Per app: privacy policy (German where offered),
store description and structured disclosure, provider website and privacy blog
posts, and advertisements encountered in use. Each document is archived with
source, language, URL, capture date, SHA-256 hash, and app version, then
segmented into coding units at heading and paragraph boundaries so that a
40-page policy and a three-line store blurb stay commensurable. Advertising is a
secondary corpus: coded where found, but its absence is not a finding, because
ad delivery cannot be reproduced.

\textbf{Step 3 --- Walkthrough.} Twelve apps, balanced 6 fitness / 6 fertility
and 3 Android / 3 iOS, inspected manually following the cognitive
walkthrough~\parencite{wharton1994cognitive} with one adaptation: the tasks are
\emph{privacy} tasks. For each, the four standard questions are asked at every
step. The tasks are: find what is collected; find what is \emph{derived}; find
what happens after a breach and how incidents are handled; reach the policy from
inside the app; stop third-party sharing; export; delete. The second and third
are expected to fail, and the walkthrough records \emph{where}.

\textbf{Step 4 --- Coding.} Three independent axes. \emph{Axis A:} harm type,
using the seven Citron and Solove categories as the backbone, plus four probe
codes (\textsc{inference}, \textsc{breach-consequence},
\textsc{incident-handling}, \textsc{concrete-consequence}), plus the
Knabenhans et al.\ pathway stage --- component, capability, or harm. \emph{Axis
B:} stance --- acknowledged, hedged, minimised, deflected, or absent (the last
asserted only at app level). \emph{Axis C:} which PET principle is invoked, and
whether protection is framed as secrecy or as harm prevention, with a separate
code for reassurance carrying no verifiable commitment. Every category gets a
definition, a positive example, and a mandatory counter-example; without the
counter-example the category has no boundary and both human and model drift.

\textbf{Coding and validation.} The corpus is too large to code entirely by
hand, so coding is LLM-assisted: one versioned prompt per axis, a pinned model
at temperature 0, three runs per unit with majority vote, structured output in
which every code carries a verbatim span --- a code without a span is discarded
automatically. \textbf{The calibration is deliberately asymmetric.} This study's
headline claims are claims of \emph{absence}, and the two error types are not
equally dangerous: a false positive inflates a count that ships with its quote
and any reader can check it, whereas a false negative manufactures the very
silence the thesis sets out to establish. So alongside the usual precision audit
and per-code Cohen's $\kappa$ against a manually double-coded 20\,\% subsample
--- reported \emph{with} prevalence, since $\kappa$ collapses on rare codes even
at near-total agreement~\parencite{cohen1960coefficient} --- there is a
\textbf{recall audit}: misses are counted on the subsample, and a lexical sweep
independent of the model searches the whole corpus for harm-adjacent vocabulary
in both languages, with every uncoded hit read by hand. An absence is therefore
never reported as \enquote{the provider does not mention X} but as
\enquote{no mention was found by an instrument whose measured recall for X is
$r$}, with $r$ given.

\textbf{Analysis.} Descriptive. With 50 apps and rare codes, no inferential
claim is made: prevalence per code at app and unit level, stance distributions
as counts, PET principle crossed with framing, and for RQ4 a pairwise comparison
of each surface's code set --- the headline being how often a code present in
the policy is absent from the screen where the user actually decides.

\textbf{Ethics.} No human subjects: public documents and the researcher's own
use of public apps. Confirmation will nonetheless be obtained in writing before
collection. Walkthroughs use throwaway accounts and synthetic data only; no real
personal data is entered into any studied app, and store terms are checked
before any automated retrieval.

\textbf{Limitations.} The study establishes that a harm type is not mentioned,
never \emph{why} --- motive is not attributed. It is a snapshot of a frozen
corpus, in one market. It says nothing about whether the communication is true;
that question belongs to the measurement literature cited above. And the
recall audit bounds, but does not eliminate, the risk that harm vocabulary
neither the codebook nor the sweep anticipated passed unseen.

%---------------------------- 4. Plan ---------------------------------
\section{Timeline and Expected Contributions}

Twelve weeks from registration, assuming registration on 2~November~2026; the
statutory three-month deadline would then fall on 2~February~2027, leaving the
final week as slack.

\begin{table}[h]
\centering\footnotesize
\begin{tabularx}{\textwidth}{@{}llX@{}}
\toprule
\textbf{Phase} & \textbf{Period} & \textbf{Deliverable} \\
\midrule
Preparation  & to 30 Oct 2026      & Proposal approved; five-app pilot; instruments v1 frozen \\
Weeks 1--4   & 2--29 Nov 2026      & Sample frozen ($n = 50$); corpus collected, hashed, segmented \\
Weeks 5--6   & 30 Nov--13 Dec 2026 & Twelve privacy walkthroughs \\
Weeks 7--8   & 14--27 Dec 2026     & LLM coding runs; manual validation subsample \\
Weeks 9--10  & 28 Dec--10 Jan 2027 & Agreement, precision and recall reports; results \\
Weeks 11--12 & 11--24 Jan 2027     & Full draft to supervisor; revisions \\
Buffer       & to 2 Feb 2027       & Final checks and submission \\
\bottomrule
\end{tabularx}
\end{table}

\noindent Introduction, background, and related work are drafted alongside
collection, so week~11 assembles rather than writes.

\textbf{Expected contributions.} (1)~The first systematic map of harm
communication across fitness and fertility trackers in the German market.
(2)~Quantified figures for four specific silences --- inference, breach
consequences, incident handling, concrete consequence --- per app and per
surface. (3)~A reusable, versioned codebook that operationalises the Citron and
Solove typology for corporate text; existing schemes for privacy policies code
data practices, this one codes harm. (4)~A privacy-focused cognitive walkthrough
instrument. (5)~A reproducible LLM-assisted pipeline reported honestly, with
agreement, precision \emph{and} recall published whatever they turn out to be.

\section{Points for Confirmation}

To keep work moving, this proposal decides every question left open after the
meeting of 4~September and the written input of 5~September. Each decision is
listed here to be confirmed or overturned; overturning one changes only the
corresponding part of Section~3.

\begin{enumerate}[itemsep=1pt,topsep=2pt]
  \item \textbf{Scope is document and interface analysis, not a user survey.}
    The 5~September input lists only document and interface sources, which
    settles the question raised in the meeting notes.
  \item \textbf{Citron and Solove as the coding backbone, Knabenhans et al.\ as
    a pathway annotation on top} --- used together rather than merged, so
    neither loses its original meaning.
  \item \textbf{\enquote{Highest rated} means $\geq 1{,}000$ ratings, then
    ranked by average rating.} Without a floor on rating count, the phrase
    selects obscure apps with a handful of five-star reviews.
  \item \textbf{50 apps is the target, not a first pass}, stratified 30/20, with
    a separate five-app pilot to test the instruments.
  \item \textbf{Walkthroughs on 12 apps, balanced 6/6 and 3/3 across platforms}
    --- balance keeps the RQ4 comparisons from being confounded with sample
    size. The alternative is walking through all 50 and shrinking the document
    sample.
  \item \textbf{A second coder} on part of the validation subsample would
    materially strengthen the design. Can one be arranged through the chair or
    MPI-SP? If not, the limitation is reported as such.
  \item \textbf{Ethics:} no human subjects are involved, so no review is
    expected to be needed --- but written confirmation is requested before
    collection starts.
\end{enumerate}

\medskip
\noindent\footnotesize
The long version of this proposal (\texttt{proposal.pdf}, in the repository)
argues every decision above in full, adds three further points for confirmation,
and carries the complete methodology from which the thesis chapter will be
written.

%---------------------------- References ------------------------------
\printbibliography[title={References}]

\end{document}
